\documentclass{beamer}

\usepackage[english,UKenglish]{babel}
\usepackage{tikz}
\usepackage{datetime}

\newcommand{\image}[2][]{%
  \IfFileExists{#2}%
    {\includegraphics[#1]{#2}}%
    {%
      \begingroup\fboxsep=-\fboxrule
      \fbox{%
        \rule{150pt}{0pt}%
        \rule{0pt}{100pt}%
      }%
      \endgroup%
    }%
}

\newdate{swat_pres_date}{16}{05}{2014}
\newdate{defence_date}{19}{06}{2014}

\definecolor{cwi-red}{cmyk}{0,.9058,.7068,.251}
\definecolor{cwi-blue}{cmyk}{.7,.4,.4,0}
\definecolor{cwi-green}{cmyk}{.41,.18,.43,.1}
\definecolor{cwi-orange}{cmyk}{.14,.35,.65,.15}
\definecolor{cwi-yellow}{cmyk}{.14,.20,.65,.16}

\usetheme{Madrid}
\useinnertheme{rectangles}
\usecolortheme[named=cwi-blue]{structure}
\setbeamertemplate{blocks}[default]
\setbeamercolor{alerted text}{fg=cwi-red}
\beamertemplatenavigationsymbolsempty

\setbeamertemplate{caption}{%
  \scriptsize%
    {\usebeamercolor[fg]{section in toc}\insertcaptionname\ \insertcaptionnumber:}%
    \ \insertcaption%
}

\newlength{\wideitemsep}
\setlength{\wideitemsep}{\itemsep}
\addtolength{\wideitemsep}{.6em}
\let\olditem\item
\renewcommand{\item}[1][\wideitemsep]{\setlength{\itemsep}{#1}\olditem}
\def\sitem{\item[.2em]}

\title[Interaction Design for Molecule Param.]{Interaction Design for Fragment-Based Molecule Parameterisation}
\author[Jimi van der Woning]{%
 Jimi M. van der Woning\\%
 \url{Jimi.vanderWoning@student.uva.nl}%
}
\institute[UvA, CWI]{%
 \parbox{4.1cm}{%
  \center%
  \scshape%
  \includegraphics[height=1.5cm]{uva_logo}\\[.1cm]%
  University of Amsterdam%
 }%
 \parbox{.5cm}{%
  \center%
  \vspace{1.5cm}%
  \ \\[.1cm]%
  and%
 }%
 \parbox{5.2cm}{%
  \center%
  \scshape%
  \includegraphics[height=1.5cm]{cwi_logo}\\[.1cm]%
  Centrum Wiskunde \& Informatica%
 }%
 \\[1cm]%
 {\bf Supervisors:} Mohammed El-Kebir, Gunnar W. Klau, and Tijs van der Storm%
}
\date{\displaydate{defence_date}}
\ddmmyyyydate




\begin{document}

\setlength{\itemsep}{20pt}



\begin{frame}
 \titlepage
\end{frame}



\addtobeamertemplate{frametitle}{\def\insertframetitle{\insertsectionhead\ \insertcontinuationcountroman}\let\insertframesubtitle\insertsubsectionhead}{%
\begin{tikzpicture}[remember picture,overlay]
\node[anchor=north east,yshift=2pt] at (current page.north east) {\includegraphics[height=0.8cm]{cwi_logo}};
\node[anchor=north east,yshift=2pt,xshift=-2cm] at (current page.north east) {\includegraphics[height=0.8cm]{uva_logo}};
\end{tikzpicture}}



\section{Introduction}
\subsection{Chemical concepts}
\begin{frame}{x}{}
  \begin{columns}
   \column{0.4\textwidth}
     \begin{figure}
      \centering
      \image[width=\textwidth]{img/molecule.pdf}
      \caption{Ethanamine molecule.}
     \end{figure}
  
   \column{0.5\textwidth}
    \begin{itemize}
      \item Everything consists of molecules;
      \item Molecules consist of atoms and bonds;
      \item Atoms have a charge.
    \end{itemize}
  \end{columns}
\end{frame}


\subsection{Problem analysis}
\begin{frame}{x}{}
  \begin{columns}
   \column{0.4\textwidth}
    \begin{itemize}
     \item<1-> Molecule simulations gain in importance;
     \item<2-> Current methods are too slow;
     \item<3-> Possible to parameterise a molecule based on fragments of other molecules;
     \item<4-> Only doable by chemists.
    \end{itemize}

   \column{0.6\textwidth}<3->
    \begin{figure}
     \vspace{-1em}
     \frame{\image[width=.42\textwidth]{../thesis/img/match_target.pdf}}%
     \qquad
     \frame{\image[width=.42\textwidth]{../thesis/img/match_none.pdf}}%
     \\[.3cm]
     \frame{\image[width=.42\textwidth]{../thesis/img/match_good.pdf}}%
     \qquad
     \frame{\image[width=.42\textwidth]{../thesis/img/match_bad.pdf}}%
     \\[.3cm]
     \vspace{-.5cm}
     \caption{Fragment matching.}
    \end{figure}
  \end{columns}
\end{frame}



\section{Research questions}
\begin{frame}{x}{}
 \begin{enumerate}
  \item<1-> Is it feasible to do manual partial charge assignment in a reasonable amount of time, i.e.\ faster than performing computations?
  \item<2-> Compared to the computed charges, how accurate are the charges obtained using manual charge assignment?
  \item<3-> How should chemists interact with a system for finding atomic partial charges using fragments of other molecules?
 \end{enumerate}
\end{frame}



\section{Approach}
\begin{frame}{x}{}
 \begin{itemize}
  \item<1-> No existing systems;
  \item<2-> Design, implement, and compare two versions;
  \item<3-> Use external tools for:
   \begin{itemize}
    \sitem Molecule visualisation: \texttt{OpenBabel}~\cite{oboyle2011open};
    \sitem Fragment finding: \texttt{mop/fragments}~\cite{elkebir2014molecule}.
   \end{itemize}
 \end{itemize}
\end{frame}



\section{Interaction Design}
\begin{frame}{x}{}
 \begin{columns}
  \column{.45\textwidth}
   \begin{block}{Reactive version}
    \begin{enumerate}
     \item<2-> Insert a molecule;
     \item<3-|alert@11> Select atom(s);
     \item<4-|alert@11> Compare fragments;
     \item<5-|alert@11> Pick a fragment;
     \item<9-> Fine-tune charges;
     \item<10-> Download result.
    \end{enumerate}
   \end{block}

  \column{.45\textwidth}
   \begin{block}{Proactive version}
    \begin{enumerate}
     \item<2-> Insert a molecule;
     \item<6-> Click ``Start'' button;
     \item<7-|alert@11> Inspect suggested fragment;
     \item<8-|alert@11> Accept/Reject fragment;
     \item<9-> Fine-tune charges;
     \item<10-> Download result.
    \end{enumerate}
   \end{block}

 \end{columns}
\end{frame}



\section{Implementation}
\begin{frame}{x}{}
 \vspace{-1.2em}
 \mbox{The Online tool for Fragment-based Molecule Parameterisation (\textsc{OFraMP})}
 \begin{center}
  Demonstration: parameterisation of propanamide~(\texttt{NC(=O)CC}).
 \end{center}
 \vspace{-.8em}
 \begin{figure}
  \image[width=.8\textwidth]{img/demo.png}
 \end{figure}
\end{frame}



\section{Evaluation}
\begin{frame}{x}{}
 \begin{itemize}
  \item<1-> On-line user study;
   \begin{itemize}
    \sitem 16 computational chemists invited, 10 participated;
    \sitem Test both versions.
   \end{itemize}
  \item<2-> Two groups:
   \begin{itemize}
    \sitem Reactive-first: 6;
    \sitem Proactive-first: 4.
   \end{itemize}
  \item<3-> For each version:
   \begin{enumerate}
    \sitem Take the guided demo;
    \sitem Parameterise a small molecule;
    \sitem Parameterise a large molecule;
    \sitem Rate using SUS questionnaire.
   \end{enumerate}
  \item<4-> General questionnaire at the end.
 \end{itemize}
\end{frame}

\section{Results}
\subsection{Time consumption}
\begin{frame}{x}{}
\begin{figure}
\image[width=\textwidth]{../thesis/img/graphs/1a_02.pdf}
\end{figure}
\end{frame}


\subsection{Result accuracy}
\begin{frame}{x}{}
\begin{figure}
\image[width=\textwidth]{../thesis/img/graphs/1a_01.pdf}
\end{figure}
\end{frame}


\subsection{Usability rating}
\begin{frame}{x}{}
 \begin{columns}
  \column{.45\textwidth}
   \begin{itemize}
    \item Reactive version: 76
     \begin{itemize}
      \sitem \emph{Above} average~\cite{sauro2011measuring};
      \sitem `Good'~\cite{bangor2009determining}.
     \end{itemize}
    \item Proactive version: 61
     \begin{itemize}
      \sitem \emph{Below} average~\cite{sauro2011measuring};
      \sitem `OK'~\cite{bangor2009determining}.
     \end{itemize}
   \end{itemize}

  \column{.45\textwidth}
   \begin{figure}
    \image[width=\textwidth]{../thesis/img/graphs/4a_10.pdf}
    \vspace{-2em}
    \caption{SUS scores.}
   \end{figure}

 \end{columns}
\end{frame}


\subsection{Comments}
\begin{frame}{x}{}
 \begin{columns}
  \column{.45\textwidth}
   \begin{itemize}
    \item<1-> \emph{All} users preferred reactive version;
    \item<2-> Interface very clear and responsive;
    \item<3-> Easy to learn;
   \end{itemize}

  \column{.45\textwidth}
   \begin{itemize}
    \item<4-> ``Excellent job, quality and functionality!''
    \item<5-> ``Nice work!''
    \item<6-> ``I think that the system has great potential!''
   \end{itemize}

 \end{columns}
 \vspace{1em}
 \begin{columns}
  \column{.45\textwidth}
   \begin{itemize}
    \item<7-> `Finished' popup confusing;
    \item<8-> Charge edit missing;
    \item<9-> Fragment matching not perfect.
   \end{itemize}

  \column{.45\textwidth}
   \ 

 \end{columns}
\end{frame}



\section{Conclusion}
\begin{frame}{x}{}
\begin{itemize}
\item<1-> Fragment-based molecule parameterisation \emph{is} possible;
\item<2-> Proactive version is quicker;
\item<3-> Reactive version yields more accurate results;
\item<4-> Reactive version was preferred by the users;
\item<5-> With a few improvements, \textsc{OFraMP} reactive will be of great value;
\item<6-> In general, non-automated systems might work better for complex, unexplored tasks.
\end{itemize}
\end{frame}



\section*{Bibliography}
\renewcommand{\item}{\olditem}
\begin{frame}[allowframebreaks]{x}{}
\nocite{badre2002shaping, blair2008user, brooke2013sus, canzar2012charge, horvitz1999principles, lewis2013umux, malde2011automated, norman2002design, sauro2011measuring, tullis2004comparison, wohlin2003empirical}
\bibliographystyle{abbrv}
\tiny{\bibliography{../thesis/references}}
\end{frame}



\end{document}
